\documentclass{exam-zh}
\usepackage{edu-practice}

\examsetup{
  question/show-answer = true,
  fillin/show-answer = true,
  solution/show-solution = show-stay,
}

% ---------- 教师版解答样式 ----------
\newtcolorbox{solutionblock}{
  enhanced, breakable,
  colback=edu-blue-bg,
  colframe=edu-blue!25,
  boxrule=0pt,
  borderline west={2pt}{0pt}{edu-blue!50},
  arc=0pt,
  left=3mm, right=3mm, top=2mm, bottom=2mm,
  before skip=2mm,
  after skip=3mm,
  fontupper=\small,
}

% 解答步骤编号
\newcounter{solstep}
\newcommand{\solstep}[2]{%
  \stepcounter{solstep}%
  \par\medskip
  \noindent\hangindent=6mm
  {\color{edu-blue}\bfseries\textcircled{\scriptsize\thesolstep}\enspace #1}\enspace #2\par
}

\begin{document}

\section*{立体几何距离与角 · 练习（教师版）}

\section*{一、选择题（每题 8 分，共 16 分）}


\begin{question}[points=8]
  直线 $l$ 与平面 $\alpha$ 斜交，$l$ 上一点 $P$（斜足外）在 $\alpha$ 内的射影为 $P'$，斜足为 $A$。下列说法正确的是（　　）。
  \begin{hintbox}
\textbf{提示：}线面角 = 斜线与其在面内射影的夹角。先有垂线才有射影。\end{hintbox}

  \begin{choices}
    \item 直线 $l$ 与平面 $\alpha$ 所成的角是 $\angle PAP'$
    \item 直线 $l$ 与平面 $\alpha$ 所成的角是 $\angle lAA'$，其中 $AA'\perp\alpha$
    \item 直线 $l$ 与平面 $\alpha$ 所成的角是 $l$ 与 $\alpha$ 内任意一条直线的夹角
    \item 线面角可以大于 $90^\circ$
  \end{choices}
  \paren[A]
  \begin{solutionblock}
    线面角定义为斜线与其在面内射影的夹角。$P$ 的射影为 $P'$，斜足为 $A$，射影即 $AP'$，故线面角为 $\angle PAP'$，选 A。B 误在再作垂线混淆对象；C 射影是唯一确定的，不是任意直线；D 线面角范围为 $[0^\circ,90^\circ]$。
  \end{solutionblock}
\end{question}



\begin{question}[points=8]
  用余弦定理算得两条异面直线平移后的夹角 $\theta$ 满足 $\cos\theta=-\dfrac{\sqrt2}{2}$。则这两条异面直线所成角为（　　）。
  \begin{hintbox}
\textbf{提示：}异面直线所成角范围是 $(0^\circ,90^\circ]$，算出钝角要取补角。\end{hintbox}

  \begin{choices}
    \item $45^\circ$
    \item $135^\circ$
    \item $60^\circ$
    \item $120^\circ$
  \end{choices}
  \paren[A]
  \begin{solutionblock}
    异面直线所成角规定为锐角或直角，范围 $(0^\circ,90^\circ]$。$\cos\theta=-\dfrac{\sqrt2}{2}$ 对应 $\theta=135^\circ$（钝角），须取补角 $180^\circ-135^\circ=45^\circ$（等价于对余弦取绝对值 $\left|-\dfrac{\sqrt2}{2}\right|=\dfrac{\sqrt2}{2}$）。故所成角为 $45^\circ$，选 A。
  \end{solutionblock}
\end{question}


\section*{二、填空题（每题 10 分，共 20 分）}


\begin{question}[points=10]
  正方体 $ABCD\text{-}A_1B_1C_1D_1$ 的棱长为 $2$。则顶点 $A_1$ 到对角面 $BDD_1B_1$ 的距离为 $\underline{\hspace{2em}}$。
  \begin{hintbox}
\textbf{提示：}点面距、垂足位置不显然时，用等体积法：选一个三棱锥列 $V$ 相等。可试 $V_{A_1\text{-}BDD_1}=V_{D\text{-}A_1BD_1}$。\end{hintbox}

  \paren[$\sqrt{2}$]
  \begin{solutionblock}
    设正方体坐标为 $A(0,0,0)$，$B(2,0,0)$，$D(0,2,0)$，$A_1(0,0,2)$，$D_1(0,2,2)$。平面 $BDD_1B_1$ 由 $B,D,D_1$ 确定，法向量可取 $(1,1,0)$，方程为 $x+y=2$。点 $A_1(0,0,2)$ 到该平面的距离为 $\dfrac{|0+0-2|}{\sqrt{1^2+1^2}}=\sqrt2$。
  \end{solutionblock}
\end{question}



\begin{question}[points=10]
  正四棱锥 $P\text{-}ABCD$ 的底面边长为 $2$，侧棱 $PA=3$，$O$ 为底面中心。则侧棱 $PA$ 与底面 $ABCD$ 所成角的正切值为 $\underline{\hspace{2em}}$。
  \begin{hintbox}
\textbf{提示：}侧棱与底面的线面角 = 斜线 $PA$ 与射影 $AO$ 的夹角。先求高 $PO$ 和 $AO$。\end{hintbox}

  \paren[$\dfrac{\sqrt{14}}{2}$]
  \begin{solutionblock}
    正四棱锥顶点在底面射影为底面中心 $O$，$PO\perp$ 底面，$AO$ 为 $PA$ 在底面射影，$\angle PAO$ 即所求线面角。$AO=\dfrac12 AC=\dfrac12\cdot 2\sqrt2=\sqrt2$，$PO=\sqrt{PA^2-AO^2}=\sqrt{9-2}=\sqrt7$。故 $\tan\angle PAO=\dfrac{PO}{AO}=\dfrac{\sqrt7}{\sqrt2}=\dfrac{\sqrt{14}}{2}$。
  \end{solutionblock}
\end{question}


\section*{三、解答题（共 60 分）}

\needspace{8\baselineskip}

\begin{problem}[points=12]
    正方体 $ABCD\text{-}A_1B_1C_1D_1$ 的棱长为 $2$，$O$ 为 $AC$ 与 $BD$ 的交点。
\begin{enumerate}[label=(\arabic*)]
  \item 证明 $BD_1\perp AC$；
  \item 求截面 $AB_1C$ 与底面 $ABCD$ 所成二面角 $A\text{-}AC\text{-}B_1$ 的余弦值。
\end{enumerate}

    \begin{hintbox}
\textbf{提示：}（1）用三垂线定理：$BD\perp AC$，$D_1D\perp$ 底面，得 $BD_1\perp AC$。（2）$BD_1\perp AC$，$BD\perp AC$，$\angle D_1BD$ 是平面角。\end{hintbox}

  \setcounter{solstep}{0}
  \begin{solutionblock}
    \textbf{答案：}(1) 见解析；(2) $\dfrac{\sqrt3}{3}$。\par\medskip
    \solstep{(1) 三垂线定理证 $BD_1\perp AC$}{正方形对角线 $BD\perp AC$，$D_1D\perp$ 底面 $ABCD$，$BD$ 为 $BD_1$ 在底面射影，由三垂线定理得 $BD_1\perp AC$。}
    \solstep{(2) 定棱与平面角顶点}{截面 $AB_1C$ 与底面 $ABCD$ 的棱为 $AC$。取 $AC$ 中点 $O$（$O$ 在棱上）。正方形中 $OB\perp AC$（对角线互相垂直平分）；$BB_1\perp$ 底面，$OB$ 为 $OB_1$ 在底面射影，由三垂线定理 $OB_1\perp AC$。故 $\angle B_1OB$ 为平面角（两边 $OB$、$OB_1$ 都 $\perp$ 棱 $AC$，顶点 $O$ 在棱上）。}
    \solstep{(2) 在直角三角形里算}{$OB=\dfrac12 BD=\dfrac12\cdot2\sqrt2=\sqrt2$；$\text{Rt}\triangle B_1OB$ 中 $B_1B=2$，$OB_1=\sqrt{B_1B^2+OB^2}=\sqrt{4+2}=\sqrt6$。}
    \solstep{(2) 解余弦}{$\cos\angle B_1OB=\dfrac{OB}{OB_1}=\dfrac{\sqrt2}{\sqrt6}=\dfrac{\sqrt3}{3}$，即所求二面角余弦值为 $\dfrac{\sqrt3}{3}$。}
  \end{solutionblock}
\end{problem}


\needspace{8\baselineskip}

\begin{problem}[points=14]
    在三棱锥 $P\text{-}ABC$ 中，$PA\perp$ 平面 $ABC$，$PA=2$，$AB=AC=2$，$\angle BAC=120^{\circ}$。求点 $A$ 到平面 $PBC$ 的距离。

    \begin{hintbox}
\textbf{提示：}点 $A$ 到面 $PBC$ 的距离用等体积法：$V_{P\text{-}ABC}=V_{A\text{-}PBC}$。先算 $S_{\triangle ABC}$ 和 $S_{\triangle PBC}$。注意 $PA\perp$ 底面，可先算 $PB$、$PC$。\end{hintbox}

  \setcounter{solstep}{0}
  \begin{solutionblock}
    \textbf{答案：}$\dfrac{2\sqrt5}{5}$\par\medskip
    \solstep{算 $S_{\triangle ABC}$}{$S_{\triangle ABC}=\dfrac12\cdot AB\cdot AC\cdot\sin120^\circ=\dfrac12\cdot2\cdot2\cdot\dfrac{\sqrt3}{2}=\sqrt3$。}
    \solstep{算 $PB$、$PC$、$BC$}{$PA\perp$ 底面，$PB=PC=\sqrt{PA^2+AB^2}=\sqrt{4+4}=2\sqrt2$；$BC=\sqrt{AB^2+AC^2-2AB\cdot AC\cos120^\circ}=\sqrt{4+4+4}=2\sqrt3$。}
    \solstep{算 $S_{\triangle PBC}$}{$\triangle PBC$ 等腰（$PB=PC=2\sqrt2$），取 $BC$ 中点 $M$，$PM=\sqrt{PB^2-\left(\dfrac{BC}{2}\right)^2}=\sqrt{8-3}=\sqrt5$，$S_{\triangle PBC}=\dfrac12\cdot BC\cdot PM=\dfrac12\cdot2\sqrt3\cdot\sqrt5=\sqrt{15}$。}
    \solstep{等体积法列方程}{设 $A$ 到面 $PBC$ 距离为 $d$。$V_{P\text{-}ABC}=\dfrac13 S_{\triangle ABC}\cdot PA=\dfrac13\cdot\sqrt3\cdot2=\dfrac{2\sqrt3}{3}$；$V_{A\text{-}PBC}=\dfrac13 S_{\triangle PBC}\cdot d=\dfrac13\cdot\sqrt{15}\cdot d$。}
    \solstep{解 $d$}{由 $\dfrac{2\sqrt3}{3}=\dfrac13\cdot\sqrt{15}\cdot d$，得 $d=\dfrac{2\sqrt3}{\sqrt{15}}=\dfrac{2}{\sqrt5}=\dfrac{2\sqrt5}{5}$。}
  \end{solutionblock}
\end{problem}


\needspace{8\baselineskip}

\begin{problem}[points=12]
    在正三棱柱 $ABC\text{-}A_1B_1C_1$ 中，底面边长为 $2$，侧棱 $AA_1=2$，$D$ 为 $BC$ 的中点。
\begin{enumerate}[label=(\arabic*)]
  \item 证明 $A_1D\perp BC$；
  \item 求直线 $A_1B$ 与底面 $ABC$ 所成角的正切值。
\end{enumerate}

    \begin{hintbox}
\textbf{提示：}（1）$AD\perp BC$（正三角形中线即高），$A_1A\perp$ 底面，由三垂线定理得 $A_1D\perp BC$。（2）$A_1B$ 与底面所成角 = $A_1B$ 与射影 $AB$ 的夹角，在 $\text{Rt}\triangle A_1AB$ 中算。\end{hintbox}

  \setcounter{solstep}{0}
  \begin{solutionblock}
    \textbf{答案：}(1) 见解析；(2) $\tan=1$（即 $45^\circ$）。\par\medskip
    \solstep{(1) 三垂线定理证 $A_1D\perp BC$}{正三角形中 $AD$ 为 $BC$ 边中线即高，$AD\perp BC$；侧棱 $AA_1\perp$ 底面 $ABC$，$AD$ 为 $A_1D$ 在底面射影，由三垂线定理得 $A_1D\perp BC$。}
    \solstep{(2) 定射影与线面角}{$AA_1\perp$ 底面，$B$ 为斜足，$AB$ 为 $A_1B$ 在底面射影，$\angle A_1BA$ 即直线 $A_1B$ 与底面所成角。}
    \solstep{(2) 在直角三角形里算}{$\text{Rt}\triangle A_1AB$ 中 $A_1A=2$（对边）、$AB=2$（邻边）。}
    \solstep{(2) 解正切}{$\tan\angle A_1BA=\dfrac{A_1A}{AB}=\dfrac{2}{2}=1$，即线面角为 $45^\circ$。}
  \end{solutionblock}
\end{problem}


\needspace{8\baselineskip}

\begin{problem}[points=10]
    正方体 $ABCD\text{-}A_1B_1C_1D_1$ 的棱长为 $2$。求点 $D_1$ 到直线 $AC$ 的距离。

    \begin{hintbox}
\textbf{提示：}连 $BD$ 交 $AC$ 于 $O$。$D_1D\perp$ 底面，$BD\perp AC$，由三垂线定理 $D_1O\perp AC$，$D_1O$ 即为所求。\end{hintbox}

  \setcounter{solstep}{0}
  \begin{solutionblock}
    \textbf{答案：}$\sqrt{6}$\par\medskip
    \solstep{找垂足 $O$}{连 $BD$ 交 $AC$ 于 $O$（正方形对角线交点），则 $O$ 为 $AC$ 中点。}
    \solstep{三垂线定理定垂线}{$D_1D\perp$ 底面，$BD\perp AC$，$DO$ 为 $D_1O$ 在底面射影，由三垂线定理得 $D_1O\perp AC$，$D_1O$ 即 $D_1$ 到 $AC$ 的距离。}
    \solstep{在直角三角形里算}{$\text{Rt}\triangle D_1DO$ 中 $D_1D=2$，$DO=\dfrac12 BD=\dfrac12\cdot2\sqrt2=\sqrt2$，故 $D_1O=\sqrt{D_1D^2+DO^2}=\sqrt{4+2}=\sqrt6$。}
  \end{solutionblock}
\end{problem}


\needspace{8\baselineskip}

\begin{problem}[points=12]
    在棱长为 $2$ 的正四面体 $P\text{-}ABC$ 中，$O$ 为底面 $\triangle ABC$ 的中心。
\begin{enumerate}[label=(\arabic*)]
  \item 求点 $P$ 到平面 $ABC$ 的距离；
  \item 设 $E$、$F$ 分别为 $PA$、$BC$ 的中点，求异面直线 $PA$ 与 $BC$ 的距离。
\end{enumerate}

    \begin{hintbox}
\textbf{提示：}（1）$PO\perp$ 底面，在 $\text{Rt}\triangle PAM$（$M$ 为 $BC$ 中点）里算 $PO$。（2）$E$、$F$ 为 $PA$、$BC$ 中点，$EF$ 为异面直线 $PA$ 与 $BC$ 的公垂线段，在 $\triangle PAM$ 里求 $EF$。\end{hintbox}

  \setcounter{solstep}{0}
  \begin{solutionblock}
    \textbf{答案：}(1) $\dfrac{2\sqrt6}{3}$；(2) $\sqrt2$。\par\medskip
    \solstep{(1) 找高 $PO$}{正四面体顶点射影为底面中心 $O$，$PO\perp$ 底面。设 $M$ 为 $BC$ 中点，$AM=\dfrac{\sqrt3}{2}\cdot2=\sqrt3$，$AO=\dfrac23 AM=\dfrac{2\sqrt3}{3}$。}
    \solstep{(1) 在直角三角形里算}{$\text{Rt}\triangle PAO$ 中 $PO=\sqrt{PA^2-AO^2}=\sqrt{4-\dfrac{12}{9}}=\sqrt{4-\dfrac43}=\sqrt{\dfrac83}=\dfrac{2\sqrt6}{3}$。}
    \solstep{(2) 证 $EF$ 为公垂线}{$AM\perp BC$（正三角形中线），$PM\perp BC$（$PM=\sqrt{PO^2+OM^2}$，$OM=\dfrac{\sqrt3}{3}$，$PM=\sqrt{\dfrac83+\dfrac13}=\sqrt3$），故 $BC\perp$ 面 $PAM$。$E$ 为 $PA$ 中点，$M(=F)$ 在面 $PAM$ 内。$\triangle PAM$ 等腰（$PM=AM=\sqrt3$），$E$ 为底 $PA$ 中点，$EM\perp PA$；又 $EM\subset$ 面 $PAM$，$BC\perp EM$。故 $EF=EM$ 为 $PA$ 与 $BC$ 的公垂线段。}
    \solstep{(2) 求 $EF$}{$\text{Rt}\triangle EAM$ 中 $AM=\sqrt3$，$AE=\dfrac12 PA=1$，$EF=EM=\sqrt{AM^2-AE^2}=\sqrt{3-1}=\sqrt2$。}
  \end{solutionblock}
\end{problem}


\clearpage
\section*{答案速查}


\begin{question}[points=0]
  参考答案：第1题 A；第2题 $\sqrt2$；第4题 $\dfrac{\sqrt{14}}{2}$；第7题 A。

  \paren[—]
\end{question}



\begin{question}[points=0]
  第3题 (1) 三垂线定理 $BD_1\perp AC$；(2) 余弦 $\dfrac{\sqrt3}{3}$（平面角 $\angle B_1OB$）。第5题 $\dfrac{2\sqrt5}{5}$。第6题 (1) 三垂线定理；(2) $\tan=1$。第8题 $\sqrt6$。第9题 (1) $\dfrac{2\sqrt6}{3}$；(2) $\sqrt2$。

  \paren[—]
\end{question}



\end{document}
